\manualchapter{Patch ideas}{sec:patches}
\subsection{Repeat a phrase with LOAD}
\begin{enumerate}
\item Load a game, start its music, and press SAVE just before an event.
\item Send a slow, regular trigger to LOAD. Each rising edge revisits that
moment; adjust the interval to choose how far the game advances.
\item Try different channel levels or route voices to separate effects.
\item Move SAVE to a new moment when you want to replace the phrase.
\end{enumerate}
An independent Rack clock is useful here. CLK's phase is not restored by
LOAD, so it is not a sample-accurate loop marker.

\subsection{Make the game control the patch}
Patch CLK to a clock divider, then use the divided outputs to trigger
percussion or advance a sequencer. Turn Clock Speed slowly to move both
the game and the downstream events. The correspondence is to the emulator's
frame counter, not to the game's internal music tempo.

\subsection{Play with gates}
Send separate gate patterns to Player 1 A and a direction. Use longer gates
for movement and shorter, separated gates for repeated button presses.
Try SAVE/LOAD to revisit the same starting situation while comparing
patterns. The result depends on the game's input handling.

\subsection{Alternate between two speeds}
With the main clock knob near its center, send a 0--5 V gate to MOD and turn
the attenuverter toward +100\%. The requested rate doubles during the high
part of the gate. Reduce the attenuverter for gentler changes, or make it
negative to slow the game during each gate.

\subsection{Add CV Genie}
Place \textbf{CV Genie (Input)} immediately to the right of RackNES. Select
the matching game and assign the memory locations you want to control.
Unassigned rows remain inactive, even with CV connected. Genie can alter game data such as position, palette, or the
Zelda sound-driver state. Its separate manual lists the maps and the
important setup limitations. RackNES accepts only the directly adjacent
input Genie; a chain of Genies does not extend the eight inputs.
